% Book 4, Year 2, chapter 21 -- Finite automata.
% Provenance (never printed): sources/book4/BRIEFS.md, chapter 21.
\chapter{Finite Automata}\label{ch:b2:finite-automata}

Is a typed string a decimal number, such as \code{-12.5e3}? A program can decide with one table and one loop:

\omcode{code/bachelor-2/21-finite-automata/listings/dfa_demo.rs}{1}{37}{A decimal number recogniser: eight states,
four classes of characters}\label{lst:b2:finite-automata:hook}

\omoutput{code/bachelor-2/21-finite-automata/outputs/dfa_demo.txt}{Nine strings and the states they
visit.}\label{out:b2:finite-automata:hook}

The program remembers nothing about the string but one number from 0 to 7, its \emph{state}, and it reads each
character once. Such a device is a \emph{finite automaton}. This chapter defines automata exactly, gives them a
power they seem to lack, guessing, shows that guessing adds nothing but costs states, sometimes exponentially many,
and builds automata for combinations of conditions. The next chapter connects them with the patterns of regular
expressions, and finds what no finite automaton can do.

\begin{recall}
The Year 1 volume: a finite-state machine has a finite set of states, a start state, and a \omterm{def:b2:finite-automata:dfa}{transition function} from a
state and an input to the next state; it may also produce outputs. Proofs by induction on the length of a \omterm{def:b2:finite-automata:word}{word}, or by
structural induction. Sets, maps, and the set of subsets of a set, which has $2^n$ elements when the set has $n$.
\end{recall}

\section{Words and languages}\label{sec:b2:finite-automata:words}

\begin{definition}[Alphabet, word]\label{def:b2:finite-automata:word}
An \emph{alphabet}\index{alphabet} is a finite nonempty set $\Sigma$, whose elements are \emph{letters}. A
\emph{word}\index{word (language)} over $\Sigma$ is a finite sequence of letters $w = a_1 a_2 \cdots a_n$; its
\emph{length} is $|w| = n$. The \emph{empty word} $\varepsilon$\index{empty word} has length 0. The set of all words
is $\Sigma^*$, and there are $|\Sigma|^n$ words of length $n$.
\end{definition}

The \emph{concatenation} of $u$ and $v$ is the \omterm{def:b2:finite-automata:word}{word} $uv$, $u$ followed by $v$: it is associative, $\varepsilon$ is
neutral, and $|uv| = |u| + |v|$. In the hook, $\Sigma$ is the set of characters, and the program reads
\code{"-12.5e3"} as a \omterm{def:b2:finite-automata:word}{word} of length 7.

\begin{definition}[Language]\label{def:b2:finite-automata:language}
A \emph{language}\index{language (formal)} over $\Sigma$ is a set of \omterm{def:b2:finite-automata:word}{words}, finite or infinite. For languages $L$
and $M$, the \emph{concatenation} is $LM = \{uv : u \in L,\ v \in M\}$; the \emph{powers} are $L^0 = \{\varepsilon\}$
and $L^{n+1} = L^n L$; and the \emph{star} of $L$ is $L^* = \bigcup_{n \ge 0} L^n$, the \omterm{def:b2:finite-automata:word}{words} made by concatenating
any number of \omterm{def:b2:finite-automata:word}{words} of $L$, none included.
\end{definition}

For example, $\{a, b\}^*$ is every \omterm{def:b2:finite-automata:word}{word} over $\{a, b\}$, $\{ab\}^* = \{\varepsilon, ab, abab, \dots\}$, and
$\{a\}\{b\}^*$, written $ab^*$, is the \omterm{def:b2:finite-automata:word}{words} made of one $a$ and then any number of $b$'s.

\begin{definition}[Prefix, suffix, factor, subword]\label{def:b2:finite-automata:factor}
Let $w = uxv$ with $u, x, v$ \omterm{def:b2:finite-automata:word}{words}. Then $u$ is a \emph{prefix}\index{prefix (word)} of $w$, $v$ a
\emph{suffix}\index{suffix (word)}, and $x$ a \emph{factor}\index{factor (word)}: a block of consecutive letters. A
\emph{subword}\index{subword} of $w$ is a \omterm{def:b2:finite-automata:word}{word} obtained by deleting letters of $w$, not necessarily consecutive.
\end{definition}

The \omterm{def:b2:finite-automata:word}{word} $abca$ has the \omterm{def:b2:finite-automata:factor}{prefixes} $\varepsilon, a, ab, abc, abca$; $bc$ is a \omterm{def:b2:finite-automata:factor}{factor}; $aa$ is a \omterm{def:b2:finite-automata:factor}{subword} but not a
\omterm{def:b2:finite-automata:factor}{factor}. A \omterm{def:b2:finite-automata:word}{word} of length $n$ has $n + 1$ \omterm{def:b2:finite-automata:factor}{prefixes} and $n + 1$ \omterm{def:b2:finite-automata:factor}{suffixes}.

\section{Deterministic automata}\label{sec:b2:finite-automata:dfa}

\begin{definition}[Deterministic finite automaton]\label{def:b2:finite-automata:dfa}
A \emph{deterministic finite automaton}\index{deterministic finite automaton} (DFA) is $A = (Q, \Sigma, \delta, q_0,
F)$: a finite set of \emph{states} $Q$, an \omterm{def:b2:finite-automata:word}{alphabet} $\Sigma$, a \emph{transition function}\index{transition function}
$\delta$, a partial map from $Q \times \Sigma$ to $Q$, a start state $q_0$, and a set $F \subseteq Q$ of
\emph{accepting states}\index{accepting state}. The \emph{extended} transition function $\delta^*$ on \omterm{def:b2:finite-automata:word}{words} is defined
by $\delta^*(q, \varepsilon) = q$ and $\delta^*(q, wa) = \delta(\delta^*(q, w), a)$, undefined as soon as one step is.
$A$ \emph{accepts} $w$ when $\delta^*(q_0, w)$ is defined and lies in $F$, and the \emph{\omterm{def:b2:finite-automata:language}{language} recognised} by $A$
is $L(A)$, the set of \omterm{def:b2:finite-automata:word}{words} it accepts.
\end{definition}

\begin{omfigure}
\begin{tikzpicture}[->, >={Stealth[length=1.8mm]}, node distance=17mm, every state/.style={minimum size=7mm,
    inner sep=1pt}, font=\scriptsize, auto]
\node[state, initial, initial text=] (0) {0};
\node[state] (1) [above right=8mm and 6mm of 0] {1};
\node[state, accepting] (2) [right=32mm of 0] {2};
\node[state] (3) [right of=2] {3};
\node[state, accepting] (4) [right of=3] {4};
\node[state] (5) [below=12mm of 3] {5};
\node[state] (6) [right=14mm of 5] {6};
\node[state, accepting] (7) [right=14mm of 6] {7};
\path (0) edge node {sign} (1) (0) edge node[below] {digit} (2) (1) edge node[pos=0.4] {digit} (2)
  (2) edge[loop above] node {digit} () (2) edge node {point} (3) (3) edge node {digit} (4)
  (4) edge[loop above] node {digit} () (2) edge node[below left] {e} (5) (4) edge node[above left] {e} (5)
  (5) edge node {sign} (6) (5) edge[bend right=30] node[below] {digit} (7) (6) edge node {digit} (7)
  (7) edge[loop above] node {digit} ();
\end{tikzpicture}
\par\smallskip{\small The automaton of the hook. A missing arrow is a missing transition: the string is rejected
there. The \omterm{def:b2:finite-automata:dfa}{accepting states} have a double circle.}
\end{omfigure}

\begin{proposition}[Reading in two parts]\label{prop:b2:finite-automata:extended}
For every state $q$ and \omterm{def:b2:finite-automata:word}{words} $u, v$: $\delta^*(q, uv) = \delta^*(\delta^*(q, u), v)$, both sides being defined or
undefined together.
\end{proposition}

\begin{proof}
By induction on $|v|$. For $v = \varepsilon$ both sides are $\delta^*(q, u)$. If $v = wa$, then $\delta^*(q, uwa) =
\delta(\delta^*(q, uw), a) = \delta(\delta^*(\delta^*(q, u), w), a) = \delta^*(\delta^*(q, u), wa)$, by the definition,
the hypothesis for $w$, and the definition again; each equality holds also when one side is undefined.
\end{proof}

This is what makes an automaton a streaming program: after reading $u$, the state $\delta^*(q_0, u)$ is all that is
needed to continue, whatever $u$ was. The \code{run_from} method of \code{automaton.rs}, the library of this
chapter's listings, is $\delta^*$ written as a \omterm{def:b2:functional-programming:fold}{fold} over the letters.

\begin{definition}[Complete automaton]\label{def:b2:finite-automata:complete}
A DFA is \emph{complete}\index{complete automaton} if $\delta(q, a)$ is defined for every state $q$ and letter $a$.
Any DFA is made complete, with the same \omterm{def:b2:finite-automata:language}{language}, by adding a non-accepting \emph{sink} state $s$ and sending every
missing transition, and every letter from $s$, to $s$.
\end{definition}

\begin{definition}[Accessible, co-accessible, trim]\label{def:b2:finite-automata:trim}
A state $q$ is \emph{accessible}\index{accessible state} if $q = \delta^*(q_0, u)$ for some \omterm{def:b2:finite-automata:word}{word} $u$, and
\emph{co-accessible}\index{co-accessible state} if $\delta^*(q, v) \in F$ for some \omterm{def:b2:finite-automata:word}{word} $v$. An automaton is
\emph{trim}\index{trim automaton} if all its states are both.
\end{definition}

\begin{proposition}[Trimming]\label{prop:b2:finite-automata:trim}
Removing from a DFA every state that is not both \omterm{def:b2:finite-automata:trim}{accessible} and \omterm{def:b2:finite-automata:trim}{co-accessible}, with the transitions that touch it,
leaves the recognised \omterm{def:b2:finite-automata:language}{language} unchanged. Both sets are computed by a graph search, in time linear in the number of
transitions.
\end{proposition}

\begin{proof}
If $w$ is accepted, every state on its run is \omterm{def:b2:finite-automata:trim}{accessible} (by a \omterm{def:b2:finite-automata:factor}{prefix} of $w$) and \omterm{def:b2:finite-automata:trim}{co-accessible} (by the rest of $w$),
so the run survives the removal. Conversely, removing states and transitions can only remove runs, so no new \omterm{def:b2:finite-automata:word}{word} is
accepted. \omterm{def:b2:finite-automata:trim}{Accessible states} are those reachable from $q_0$ in the graph of the transitions; \omterm{def:b2:finite-automata:trim}{co-accessible states} are
those from which $F$ is reachable, that is, those reachable from $F$ in the reversed graph.
\end{proof}

\begin{definition}[Recognisable language]\label{def:b2:finite-automata:recognisable}
A \omterm{def:b2:finite-automata:language}{language} is \emph{recognisable}\index{recognisable language} if some finite automaton recognises it.
\end{definition}

\section{Nondeterministic automata}\label{sec:b2:finite-automata:nfa}

Consider the \omterm{def:b2:finite-automata:word}{words} over $\{a, b\}$ whose third letter from the end is $a$. Reading from the left, a program cannot
know which letter will be third from the end. It would be easy if it could \emph{guess}: on each $a$, guess ``this
is the one'', then check that exactly two letters follow.

\begin{definition}[Nondeterministic finite automaton]\label{def:b2:finite-automata:nfa}
A \emph{nondeterministic finite automaton}\index{nondeterministic finite automaton} (NFA) is $A = (Q, \Sigma, \Delta,
I, F)$ where $\Delta \subseteq Q \times \Sigma \times Q$ is a set of transitions, so that a state may have several
transitions, or none, on the same letter, and $I \subseteq Q$ is a set of start states. A \emph{run} on $w = a_1
\cdots a_n$ is a sequence $q_0 q_1 \cdots q_n$ with $q_0 \in I$ and $(q_{i-1}, a_i, q_i) \in \Delta$; $A$
\emph{accepts} $w$ if \emph{some} run on $w$ ends in $F$.
\end{definition}

\begin{omfigure}
\begin{tikzpicture}[->, >={Stealth[length=1.8mm]}, node distance=18mm, every state/.style={minimum size=7mm,
    inner sep=1pt}, font=\scriptsize, auto]
\node[state, initial, initial text=] (0) {0};
\node[state] (1) [right of=0] {1}; \node[state] (2) [right of=1] {2};
\node[state, accepting] (3) [right of=2] {3};
\path (0) edge[loop above] node {$a, b$} () (0) edge node {$a$} (1) (1) edge node {$a, b$} (2)
  (2) edge node {$a, b$} (3);
\begin{scope}[xshift=78mm]
\node[state, initial, initial text=] (p) {0}; \node[state] (q) [right of=p] {1};
\node[state, accepting] (r) [right of=q] {2};
\path (p) edge[loop above] node {$a$} () (q) edge[loop above] node {$b$} () (r) edge[loop above] node {$c$} ()
  (p) edge[dashed] node {$\varepsilon$} (q) (q) edge[dashed] node {$\varepsilon$} (r);
\end{scope}
\end{tikzpicture}
\par\smallskip{\small Left: an NFA for ``the third letter from the end is $a$''; state 0 has two transitions on $a$.
Right: an NFA for $a^*b^*c^*$ with \omterm{def:b2:finite-automata:epsilon}{spontaneous transitions} (dashed).}
\end{omfigure}

A nondeterministic automaton is not a program that guesses well; it is a definition by existence of a run. It is
also convenient to allow moves that read nothing.

\begin{definition}[Spontaneous transition, closure]\label{def:b2:finite-automata:epsilon}
An NFA \emph{with spontaneous transitions} also has a set of pairs $(p, q)$, written $p \xrightarrow{\varepsilon}
q$\index{spontaneous transition}, which a run may take without reading a letter. The
\emph{$\varepsilon$-closure}\index{epsilon-closure@$\varepsilon$-closure} of a set $S$ of states is the set of states
reachable from $S$ by spontaneous transitions alone, $S$ included: the states reachable from $S$ in the graph of
those transitions.
\end{definition}

\omcode{code/bachelor-2/21-finite-automata/listings/automaton.rs}{106}{124}{The closure by a graph search, and one
letter followed by the closure}\label{lst:b2:finite-automata:closure}

\begin{proposition}[Removing spontaneous transitions]\label{prop:b2:finite-automata:epsilon}
Let $\mathrm{cl}(q)$ be the closure of $\{q\}$. The NFA with the same states and start states, the transitions $(q, a,
r)$ for every $r$ in the closure of $\{r' : (p, a, r') \in \Delta,\ p \in \mathrm{cl}(q)\}$, no \omterm{def:b2:finite-automata:epsilon}{spontaneous
transition}, and the \omterm{def:b2:finite-automata:dfa}{accepting states} $\{q : \mathrm{cl}(q) \cap F \ne \emptyset\}$, recognises the same \omterm{def:b2:finite-automata:language}{language}.
\end{proposition}

\begin{proof}
A run of the original automaton on $a_1 \cdots a_n$ is a sequence of spontaneous moves, a letter, spontaneous moves,
and so on. Cut it just before each letter: from state $q_{i-1}$ at the cut, it moves spontaneously within
$\mathrm{cl}(q_{i-1})$, reads $a_i$, then moves spontaneously to the next cut $q_i$, which is exactly a new transition
$(q_{i-1}, a_i, q_i)$; after the last letter it moves spontaneously to $F$, which makes $q_n$ accepting in the new
automaton. Conversely, each new transition and each new \omterm{def:b2:finite-automata:dfa}{accepting state} unfolds into such moves of the original
one.
\end{proof}

\omoutput{code/bachelor-2/21-finite-automata/outputs/epsilon.txt}{The closures of $a^*b^*c^*$'s states, the automaton
without \omterm{def:b2:finite-automata:epsilon}{spontaneous transitions}, and a check on every \omterm{def:b2:finite-automata:word}{word} up to length 6.}\label{out:b2:finite-automata:epsilon}

\section{The subset construction}\label{sec:b2:finite-automata:subset}

An NFA can be run by a program that keeps the \emph{set} of states it might be in. That set takes finitely many
values, which makes it the state of a DFA.

\begin{method}[The subset construction]\label{met:b2:finite-automata:subset}
Given an NFA $(Q, \Sigma, \Delta, I, F)$, build a DFA whose states are sets of NFA states. The start state is
$\mathrm{cl}(I)$. From a set $S$ and a letter $a$, the transition goes to $\mathrm{cl}(\{r : (p, a, r) \in \Delta,\
p \in S\})$. A set is accepting if it contains a state of $F$. Build only the sets reachable from the start, with a
queue.\index{subset construction}
\end{method}

\omcode{code/bachelor-2/21-finite-automata/listings/automaton.rs}{138}{159}{The \omterm{met:b2:finite-automata:subset}{subset construction}: each new set
gets a number and joins the queue}\label{lst:b2:finite-automata:subset}

\begin{theorem}[The subset construction is correct]\label{thm:b2:finite-automata:subset-correct}
For every \omterm{def:b2:finite-automata:word}{word} $w$, the state of the DFA after reading $w$ is the set of NFA states at which some run on $w$ can end.
Hence the DFA recognises the \omterm{def:b2:finite-automata:language}{language} of the NFA, and a \omterm{def:b2:finite-automata:language}{language} is \omterm{def:b2:finite-automata:recognisable}{recognisable} by an NFA if and only if it is
\omterm{def:b2:finite-automata:recognisable}{recognisable} by a DFA.
\end{theorem}

\begin{proof}
By induction on $|w|$. For $w = \varepsilon$, the runs on $\varepsilon$ are the spontaneous moves from $I$, which end
exactly in $\mathrm{cl}(I)$, the start state. Suppose the claim holds for $w$, with set $S$. A run on $wa$ is a run on
$w$, ending in some $p \in S$, then a transition $(p, a, r)$, then spontaneous moves: it ends exactly in the closure
of the $r$'s, which is the transition of the DFA from $S$ on $a$. So the DFA's state after $w$ contains a state of $F$
if and only if some run on $w$ ends in $F$. Conversely a DFA is an NFA with one start state and at most one
transition per state and letter.
\end{proof}

\omoutput{code/bachelor-2/21-finite-automata/outputs/subset.txt}{The determinisation of ``the third letter from the
end is $a$'': eight sets of states.}\label{out:b2:finite-automata:subset}

The DFA remembers which of the last three letters were $a$: the set $\{0, 2, 3\}$ means the last letter was $b$ and
the two before it were $a$. That needs $2^3$ states, and no DFA can do with fewer.

\begin{theorem}[The blow-up is unavoidable]\label{thm:b2:finite-automata:blowup}
Let $L_k$ be the \omterm{def:b2:finite-automata:word}{words} over $\{a, b\}$ whose $k$-th letter from the end is $a$. An NFA with $k + 1$ states recognises
$L_k$, and every DFA that recognises $L_k$ has at least $2^k$ states; the \omterm{met:b2:finite-automata:subset}{subset construction} produces exactly
$2^k$.
\end{theorem}

\begin{proof}
Let $u \ne v$ be two \omterm{def:b2:finite-automata:word}{words} of length $k$, and $i$ the first position, counting from 0, where they differ; say $u$ has
$a$ there and $v$ has $b$. Append $b^i$: in $ub^i$ and $vb^i$, that position is now $k$-th from the end, so $ub^i
\in L_k$ and $vb^i \notin L_k$. In a DFA recognising $L_k$, if $\delta^*(q_0, u) = \delta^*(q_0, v)$ then by
\cref{prop:b2:finite-automata:extended} $\delta^*(q_0, ub^i) = \delta^*(q_0, vb^i)$, and the two \omterm{def:b2:finite-automata:word}{words} would be both
accepted or both rejected. So the $2^k$ \omterm{def:b2:finite-automata:word}{words} of length $k$ lead to $2^k$ different states. The NFA of the figure has
$k + 1$ states; the reachable sets of the \omterm{met:b2:finite-automata:subset}{subset construction} are $\{0\} \cup T$ for the $2^k$ sets $T \subseteq \{1,
\dots, k\}$ of positions among the last $k$ letters that hold an $a$, all reached.
\end{proof}

\omoutput{code/bachelor-2/21-finite-automata/outputs/blowup.txt}{For $k = 1$ to 12, the states of both automata, and
the proof's argument checked on every pair of \omterm{def:b2:finite-automata:word}{words} of length $k$.}\label{out:b2:finite-automata:blowup}

A set of \omterm{def:b2:finite-automata:word}{words} that must lead to pairwise different states, as here, is a \emph{fooling set}; it gives a lower bound on
the size of every DFA for the \omterm{def:b2:finite-automata:language}{language} (\cref{exo:b2:finite-automata:10}).

\begin{history}
% ledger: b2:kleene-1956
% ledger: b2:rabin-scott-1959
Finite automata come from the 1950s. Stephen Kleene's paper of 1956, in the collection \emph{Automata Studies},
started from models of nerve nets and characterised the ``events'' a finite automaton can recognise; that result is
the next chapter's main theorem. In 1959 Michael Rabin and Dana Scott defined nondeterministic automata and proved
that the \omterm{met:b2:finite-automata:subset}{subset construction} removes the nondeterminism, the theorem of this section.
\end{history}

\section{Combining automata}\label{sec:b2:finite-automata:closure}

\begin{definition}[Product automaton]\label{def:b2:finite-automata:product}
The \emph{product}\index{product automaton} of two \omterm{def:b2:finite-automata:complete}{complete} DFAs $(P, \Sigma, \delta, p_0, F)$ and $(Q, \Sigma,
\gamma, q_0, G)$ has the states $P \times Q$, the start $(p_0, q_0)$ and the transitions $(p, q) \xrightarrow{a}
(\delta(p, a), \gamma(q, a))$: it runs both automata at once.
\end{definition}

\begin{proposition}[Closure under the Boolean operations]\label{prop:b2:finite-automata:closure}
If $L$ and $M$ are \omterm{def:b2:finite-automata:recognisable}{recognisable}, so are $\Sigma^* \setminus L$, $L \cup M$ and $L \cap M$.
\end{proposition}

\begin{proof}
Take \omterm{def:b2:finite-automata:complete}{complete} DFAs for $L$ and $M$. Swapping accepting and other states of the first gives the complement: the run on
every \omterm{def:b2:finite-automata:word}{word} is defined and ends in exactly one state, accepting in one automaton if and only if not in the other. By
induction on $|w|$, the \omterm{def:b2:finite-automata:product}{product}'s state after $w$ is the pair of the two states after $w$; so with \omterm{def:b2:finite-automata:dfa}{accepting states} $F
\times G$ it recognises $L \cap M$, and with $(F \times Q) \cup (P \times G)$ it recognises $L \cup M$.
\end{proof}

Completeness matters for the complement: an incomplete automaton rejects some \omterm{def:b2:finite-automata:word}{words} by having no run at all, and
swapping its \omterm{def:b2:finite-automata:dfa}{accepting states} leaves them rejected (\cref{exo:b2:finite-automata:2}). The \omterm{def:b2:finite-automata:product}{product} of ``even length''
and ``starts with $a$'' has six states, and not all of them are useful:

\omoutput{code/bachelor-2/21-finite-automata/outputs/product.txt}{The \omterm{def:b2:finite-automata:product}{product} for the intersection, its useless
states, and the number of accepted \omterm{def:b2:finite-automata:word}{words} of each length.}\label{out:b2:finite-automata:product}

\begin{omfigure}
\begin{tikzpicture}[->, >={Stealth[length=1.8mm]}, node distance=20mm, every state/.style={minimum size=9mm,
    inner sep=0pt}, font=\scriptsize, auto, g/.style={draw=gray!60, text=gray!70}]
\node[state, initial, initial text=] (00) {0,0};
\node[state] (11) [right=24mm of 00] {1,1};
\node[state, accepting] (01) [right=24mm of 11] {0,1};
\node[state, g] (12) [below=13mm of 11] {1,2};
\node[state, g] (02) [right=24mm of 12] {0,2};
\node[state, g] (10) [below=13mm of 00] {1,0};
\path (00) edge node {$a$} (11) (00) edge node[below left] {$b$} (12)
  (11) edge[bend left=15] node {$a, b$} (01) (01) edge[bend left=15] node {$a, b$} (11)
  (12) edge[bend left=15, gray] node {$a, b$} (02) (02) edge[bend left=15, gray] node {$a, b$} (12)
  (10) edge[gray, bend right=28] node[below] {$b$} (02) (10) edge[gray, bend right=50] node[below, pos=0.3] {$a$} (01);
\end{tikzpicture}
\par\smallskip{\small The \omterm{def:b2:finite-automata:product}{product} of ``even length'' (first component) and ``starts with $a$'' (0: start, 1: began
with $a$, 2: began with $b$). Grey: $(1, 0)$ is not \omterm{def:b2:finite-automata:trim}{accessible}; $(1, 2)$ and $(0, 2)$ cannot reach the \omterm{def:b2:finite-automata:dfa}{accepting
state}. The \omterm{def:b2:finite-automata:trim}{trim automaton} keeps three states.}
\end{omfigure}

The counts, $2^{n-1}$ \omterm{def:b2:finite-automata:word}{words} of every even length $n \ge 2$, are computed from the automaton by dynamic programming:
the number of \omterm{def:b2:finite-automata:word}{words} leading to each state, letter after letter (\cref{exo:b2:finite-automata:4}).

\section{Automata as programs}\label{sec:b2:finite-automata:programs}

A DFA runs in constant time per letter: one table lookup, as in the hook. An NFA can be run either by building the
DFA first, or by simulating it on the fly, keeping the set of possible states, which costs up to the number of NFA
states per letter but builds nothing. With at most 64 states, the set fits in the bits of one machine word, and
a step is a few shifts and ORs (\cref{exo:b2:finite-automata:11}). On $L_k$, over two million letters:

\omoutput{code/bachelor-2/21-finite-automata/outputs/simulate_vs_determinise.txt}{Simulating $L_k$'s NFA with sets as
ordered trees or as bits, or building its DFA and running it, on the laptop.}\label{out:b2:finite-automata:measured}

The DFA wins while it is small. At $k = 16$ it has 65\,536 states, takes 80 ms to build, and runs slower than the
simulation by bits, since its table no longer fits in the processor's caches and each letter costs a memory access
(\cref{ch:b2:caches}). Programs that search text with patterns face the same choice, and many build the DFA lazily,
one state the first time it is needed.

\section{Exercises}\label{sec:b2:finite-automata:exercises}

\begin{exercise}[$\star$]\label{exo:b2:finite-automata:1}
Give the run of the automaton of \cref{lst:b2:finite-automata:hook} on \code{0.0}, \code{-.5}, \code{1E+10},
\code{7e3.5} and \code{+}, and say which are accepted.
\end{exercise}

\begin{exercise}[$\star$]\label{exo:b2:finite-automata:2}
The DFA with states 0 and 1, start 0, accepting 1, and the transitions $0 \xrightarrow{a} 1$ and $1 \xrightarrow{b}
1$ recognises $ab^*$. Swap its accepting and other states: which \omterm{def:b2:finite-automata:language}{language} do you get? Which \omterm{def:b2:finite-automata:word}{words} are missing from
the complement, and how is it repaired?
\end{exercise}

\begin{exercise}[$\star$]\label{exo:b2:finite-automata:3}
Build a DFA with three states that recognises the binary writings of the multiples of 3 (the \omterm{def:b2:finite-automata:word}{empty word} standing for
0), and prove it correct.
\end{exercise}

\begin{exercise}[$\star$]\label{exo:b2:finite-automata:4}
Build a \omterm{def:b2:finite-automata:complete}{complete} DFA for the \omterm{def:b2:finite-automata:word}{words} over $\{a, b\}$ that do not contain the \omterm{def:b2:finite-automata:factor}{factor} $bb$, and find a recurrence for the
number $c_n$ of such \omterm{def:b2:finite-automata:word}{words} of length $n$.
\end{exercise}

\begin{exercise}[$\star\star$]\label{exo:b2:finite-automata:5}
Apply the \omterm{met:b2:finite-automata:subset}{subset construction} to the NFA with states 0, 1, 2, start 0, accepting 2, and the transitions $0
\xrightarrow{a, b} 0$, $0 \xrightarrow{a} 1$ and $1 \xrightarrow{b} 2$. Which \omterm{def:b2:finite-automata:language}{language} is it?
\end{exercise}

\begin{exercise}[$\star\star$]\label{exo:b2:finite-automata:6}
A program computes the closure of a set by adding, once, the targets of the \omterm{def:b2:finite-automata:epsilon}{spontaneous transitions} of its states. On
the automaton of $a^*b^*c^*$, which \omterm{def:b2:finite-automata:word}{words} does it get wrong, and why?
\end{exercise}

\begin{exercise}[$\star\star$]\label{exo:b2:finite-automata:7}
Build a DFA with four states for the \omterm{def:b2:finite-automata:word}{words} over $\{a, b\}$ that contain the \omterm{def:b2:finite-automata:factor}{factor} $aba$, and explain what each state
remembers.
\end{exercise}

\begin{exercise}[$\star\star$]\label{exo:b2:finite-automata:8}
The \emph{reverse} of $w = a_1 \cdots a_n$ is $a_n \cdots a_1$. Show that if $L$ is \omterm{def:b2:finite-automata:recognisable}{recognisable}, the set of the
reverses of its \omterm{def:b2:finite-automata:word}{words} is \omterm{def:b2:finite-automata:recognisable}{recognisable}.
\end{exercise}

\begin{exercise}[$\star\star$]\label{exo:b2:finite-automata:9}
For integers $n \ge 1$ and $b \ge 2$, build a DFA with $n$ states that recognises the writings in base $b$ of the
multiples of $n$, and test it.
\end{exercise}

\begin{exercise}[$\star\star\star$]\label{exo:b2:finite-automata:10}
A \emph{fooling set} for $L$ is a set $S$ of \omterm{def:b2:finite-automata:word}{words} such that for any two $u \ne v$ in $S$ there is a \omterm{def:b2:finite-automata:word}{word} $z$ with
exactly one of $uz$ and $vz$ in $L$. Show that every \omterm{def:b2:finite-automata:complete}{complete} DFA recognising $L$ (one whose \omterm{def:b2:finite-automata:dfa}{transition function} is defined everywhere) has at
least $|S|$ \omterm{def:b2:finite-automata:trim}{accessible states}, and a partial one at least $|S| - 1$. Find a
fooling set of size 3 for the \omterm{def:b2:finite-automata:language}{language} of \cref{exo:b2:finite-automata:3}.
\end{exercise}

\begin{exercise}[$\star\star\star$]\label{exo:b2:finite-automata:11}
Write an NFA simulation that keeps the set of states as the bits of a 64-bit word, and compare its speed with the
determinised automaton for $L_k$.
\end{exercise}

\begin{exercise}[$\star\star\star$]\label{exo:b2:finite-automata:12}
Show that if $L$ is \omterm{def:b2:finite-automata:recognisable}{recognisable}, the set of the \omterm{def:b2:finite-automata:factor}{prefixes} of its \omterm{def:b2:finite-automata:word}{words} is \omterm{def:b2:finite-automata:recognisable}{recognisable}.
\end{exercise}

\section{Problem: Decoding the lighthouse's light signals}\label{sec:b2:finite-automata:problem}

\begin{problem}[{Weekend problem --- Decoding the lighthouse's light signals}]\label{pb:b2:finite-automata:1}
The lighthouse sends messages as flashes, short ($s$) or long ($l$). A message is a \omterm{def:b2:finite-automata:word}{word} over $\{s, l\}$; the harbour
reads it as it comes.

\needspace{4\baselineskip}\medskip\noindent\textbf{Part I --- Signals as words.}\par\nopagebreak
\begin{enumerate}
\item How many messages of length 10 are there, and how many of length at most 10?
\item List the \omterm{def:b2:finite-automata:factor}{prefixes}, \omterm{def:b2:finite-automata:factor}{suffixes} and \omterm{def:b2:finite-automata:factor}{factors} of length 2 of $sllsl$.
\item Is $ss$ a \omterm{def:b2:finite-automata:factor}{subword} of $slsl$? A \omterm{def:b2:finite-automata:factor}{factor}?
\item Describe $\{s, l\}^* l \{s, l\}^*$ and $\{ss, l\}^*$ in \omterm{def:b2:finite-automata:word}{words}.
\item Show that $L^* L^* = L^*$ for every \omterm{def:b2:finite-automata:language}{language} $L$.
\item Is $\{s\}^* \{l\}^* = \{s, l\}^*$? Give a \omterm{def:b2:finite-automata:word}{word} that tells.
\end{enumerate}

\needspace{4\baselineskip}\medskip\noindent\textbf{Part II --- Valid messages.}\par\nopagebreak
A message is \emph{valid} if it has an even number of long flashes and never three short flashes in a row.
\begin{enumerate}[resume, beginpenalty=10000]
\item Build a \omterm{def:b2:finite-automata:complete}{complete} DFA $P$ with two states for ``an even number of $l$''.
\item Build a \omterm{def:b2:finite-automata:complete}{complete} DFA $N$ with four states for ``no \omterm{def:b2:finite-automata:factor}{factor} $sss$''.
\item How many states has their \omterm{def:b2:finite-automata:product}{product}, and which are accepting?
\item Which states of the \omterm{def:b2:finite-automata:product}{product} are not \omterm{def:b2:finite-automata:trim}{co-accessible}? How many states does the \omterm{def:b2:finite-automata:trim}{trim automaton} keep?
\item Run it on $slslss$ and $lssslsl$.
\item Write the recurrence that counts valid messages, state by state, and compute the counts for lengths 0 to 6.
\end{enumerate}

\needspace{4\baselineskip}\medskip\noindent\textbf{Part III --- A pattern near the end.}\par\nopagebreak
The harbour must react when the fourth flash from the end is long.
\begin{enumerate}[resume, beginpenalty=10000]
\item Draw an NFA with 5 states for it.
\item What does its \omterm{met:b2:finite-automata:subset}{subset construction} remember after each flash?
\item How many states does the DFA have, and what is its start state?
\item Give two messages of length 4 that lead to different DFA states, and a \omterm{def:b2:finite-automata:factor}{suffix} that tells them apart.
\item Is the NFA's on-the-fly simulation cheaper than the DFA here? Use \cref{out:b2:finite-automata:measured}.
\item Why can the harbour not use an automaton that reads the message backwards?
\end{enumerate}

\needspace{4\baselineskip}\medskip\noindent\textbf{Part IV --- The blow-up.}\par\nopagebreak
\begin{enumerate}[resume, beginpenalty=10000]
\item For ``the $k$-th flash from the end is long'', how many states do the NFA and the DFA need?
\item Which $2^k$ \omterm{def:b2:finite-automata:word}{words} form a fooling set, and why does that prove the bound?
\item Could a different NFA with fewer than $k + 1$ states recognise the \omterm{def:b2:finite-automata:language}{language}?
\item To check also that the message is valid, the harbour runs the \omterm{def:b2:finite-automata:product}{product} of the DFA of question 19 with the \omterm{def:b2:finite-automata:trim}{trim
automaton} of question 10. At most how many states does it have for $k = 4$?
\item How does \cref{out:b2:finite-automata:blowup} check the proof for $k \le 12$?
\item State the result: the $2^k$ law, exactly $2^k$ states, none removable.
\end{enumerate}
\end{problem}
